\begin{helper}
For every natural number \(p\), the number of ordered pairs
\[
  (x,y)\in \mathbb F_2^p\times\mathbb F_2^p
\]
with \(\langle x,y\rangle=1\) is
\[
  (2^p-1)\,2^{p-1}.
\]
\end{helper}
\begin{proof}
Let \(V=\mathbb F_2^p\). If \(\langle x,y\rangle=1\), then \(x\ne0\), since the
zero vector has dot product \(0\) with every \(y\). Thus the set being counted is
the disjoint union, over all nonzero \(x\in V\), of the fibers
\[
  H_x=\{y\in V:\langle x,y\rangle=1\}.
\]
There are \(2^p-1\) choices for nonzero \(x\).

Fix such an \(x\). Choose a coordinate \(i\) with \(x_i\ne0\). Since we are over
\(\mathbb F_2\), \(x_i\) is invertible. Projection away from the \(i\)-th
coordinate gives a bijection from \(H_x\) to the vector space of functions on
the remaining \(p-1\) coordinates: given arbitrary values off \(i\), the equation
\[
  \sum_j x_jy_j=1
\]
determines the missing coordinate uniquely as
\[
  y_i=\frac{1-\sum_{j\ne i}x_jy_j}{x_i}.
\]
Conversely, every \(y\in H_x\) is recovered from its off-\(i\) coordinates by
this formula. Hence \(|H_x|=2^{p-1}\).

Summing the constant fiber size \(2^{p-1}\) over the \(2^p-1\) nonzero choices
of \(x\) gives \((2^p-1)2^{p-1}\), as required.
\end{proof}
